% !PRINTED 333
\weberpage{333}
%<<CONT

\noindent 3.\ \gesperrt{A ternary cubic form $f(x_1,x_2,x_3)$ whose discriminant is different from zero can be brought, in four different ways, into the canonical form}
\[
\varphi(y_1,y_2,y_3) = h(y_1^3+y_2^3+y_3^3) + 6\,m\,y_1y_2y_3
\]
\gesperrt{by transformation.}

Since on each of the lines $y_1=0$, $y_2=0$, $y_3=0$ there lie three inflexion points, the problem of the inflexion points is essentially identical with that of the transformation of the cubic form $f$ to the canonical form, with which we shall occupy ourselves first $^1$).

\begin{center}
\S.\,90.

\gesperrt{Transformation of the cubic form to the canonical form.}
\end{center}

In order to carry out the transformation of the ternary cubic form to the canonical form, that is, to reduce it to the solution of certain algebraic equations, we must first form the covariants for the canonical form.

Accordingly, let
\[
\varphi(y) = h(y_1^3+y_2^3+y_3^3) + 6\,m\,y_1y_2y_3,
\]
for which, as an abbreviation, we also set:
\weq{1}{\varphi(y) = h\,\Sigma y_1^3 + 6\,m\,y_1y_2y_3,}
where by the symbol $\Sigma$ we understand a sum over three terms formed from the first by cyclic permutation of the variables $y_1,y_2,y_3$.

To simplify the formulae we shall multiply the \gesperrt{Hesse} covariant by a suitable numerical factor, and accordingly
\weq{2}{\Delta(x) = \frac{1}{6^3}\left|\begin{array}{ccc} f_{11}, & f_{12}, & f_{13} \\ f_{21}, & f_{22}, & f_{23} \\ f_{31}, & f_{32}, & f_{33} \end{array}\right|}

\medskip
\noindent\rule{4cm}{0.4pt}

\smallskip
\noindent $^1$) Of the literature on the points of inflection of curves of the third order and on the theory of ternary cubic forms we mention, besides the older works of \gesperrt{Newton} and \gesperrt{Mac Laurin}: \gesperrt{Hesse}, ``Ueber die Elimination etc.'' and ``Ueber die Wendepunkte der Curven dritter Ordnung'', \gesperrt{Crelle's} Journal, Bd. 28 (1844). \gesperrt{Aronhold}, Theorie der homogenen Functionen 3\textsuperscript{ten} Grades von drei Veränderlichen, \gesperrt{Crelle's} Journal, Bd. 55 (1858). \gesperrt{Salmon}, Higher plane curves (deutsch von \gesperrt{Fiedler}).
%CONT>>
